% ============================================================
% glossary.tex —— 术语表（Glossary）
% 底本：LNM 2386，PDF p477–p484（印刷 pp.437–444）
% 体例：两列 longtable（英文术语 | 中文译名＋一句话释义）
% 说明：原书每条按 (Physics)/(Mathematics) 分两段详述；
%       本译注本按 Lead 体例压成「译名＋一句话释义」，
%       完整论述见正文对应章节（条目末给出原文所标章节）。
% ============================================================

% [原文 p.437（印刷 p.437）]
\chapter*{术语表}
\addcontentsline{toc}{chapter}{术语表}

% [Tutor Note: 原书 Glossary 每条都分 (Physics) 与 (Mathematics) 两段（部分条目只有一段），
% 分别给出物理学与数学的界定。本表按项目体例每条只保留「译名＋一句话释义」，并把原文
% 所标的章节出处附在句末；需要完整表述时请回查正文相应章节。]

\small
\begin{longtable}{@{}p{0.24\linewidth}p{0.72\linewidth}@{}}
BPS state & \textbf{BPS 态}：饱和 Bogomol'nyi 界 $M=|Z|$ 的态（$M$ 为不变量、$Z$ 为中心荷），等价地是超对称代数的短表示，因而对量子修正稳定（见第 17、20 章）。 \\
Character variety & \textbf{特征簇}：基本群 $\pi_1(M)$ 到李群 $G$ 的表示模去 $G$ 共轭作用的同构类空间 $\mathcal{R}(\Gamma,G)$，与 $M$ 上的几何结构、$G$-主丛的平坦联络及 Higgs 丛密切相关（见 9.1 节与 3.4 节）。 \\
% [原文 p.438（印刷 p.438）]
Class S theory & \textbf{S 类理论}：由六维 $\mathcal{N}=(2,0)$ 超共形场论在黎曼曲面 $C$ 上紧化所得的 4d $\mathcal{N}=2$ 量子场论，其 Seiberg–Witten 可积系统是 Hitchin 系统（见第 21 章）。 \\
Cluster variety & \textbf{簇代数簇}：把若干份 $(\C^*)^n$ 沿称为变异的特殊有理映射粘合而成的代数簇，与箭图的变异类对应；分 A 型（带预辛结构）与 X 型（带典范泊松结构）两类（见 5.5 节）。 \\
Component of maximal representations & \textbf{极大表示分支}：Toledo 数取到上界（由曲面欧拉示性数与群秩决定）的表示即极大表示，它们构成的分支与 Hitchin 分支同为高阶泰希米勒空间的基本例子（见第 11 章）。 \\
Conformal group & \textbf{共形群}：保持（伪）黎曼度量共形类的微分同胚群；例如欧氏 $n$ 球面的共形群为 $\SO(1,n+1)$，洛伦兹 $n$ 维球面的共形群为 $\SO(2,n)$（见 17.1.2 节）。 \\
% [原文 p.439（印刷 p.439）]
Coulomb branch & \textbf{库仑分支}：超对称真空模空间中低能有效理论为阿贝尔规范理论的分支；4d $\mathcal{N}=2$ 时它是特殊凯勒流形，在 S 类理论中它是黎曼曲面 $C$ 上 Hitchin 系统的底（见第 18 章）。 \\
Covering space & \textbf{覆叠空间}：连续映射 $p:\tilde X\to X$，使每个 $x\in X$ 的原像 $p^{-1}(x)$ 为若干份 $X$ 的离散无交并；允许分歧点者称为分歧覆叠（见第 2 章）。 \\
Electric-magnetic duality & \textbf{电磁对偶}：Maxwell 方程在电场与磁场互换下的不变性；量子层面需同时把电耦合取倒数 $e\mapsto e^{-1}$，在规范理论中推广为 Montonen–Olive 对偶（S 对偶），计入 $\Theta$ 角后对偶群增强为 $\SL_2(\Z)$（见第 19 章）。 \\
% [原文 p.440（印刷 p.440）]
Gauge theory & \textbf{规范理论}：以李群 $G$ 为规范群、带有局域冗余（规范对称性）的场论，如 Maxwell、Chern–Simons 与 Yang–Mills 理论；几何上由 $G$-主丛及其联络描述（见第 3 章）。 \\
Higgs branch & \textbf{希格斯分支}：超对称真空模空间中低能时规范群被完全破缺的分支；4d $\mathcal{N}=2$ 时为超凯勒流形（见第 18 章）。 \\
Higher-rank Teichmüller space & \textbf{高阶泰希米勒空间}：$G$-特征簇 $\Hom(\pi_1(S),G)/G$ 中仅由离散忠实表示构成的那些连通分支之并（$\rank(G)>1$），典型例子是 Hitchin 分支与极大表示分支（第 II 部分的主角）。 \\
% [原文 p.441（印刷 p.441）]
Hitchin component & \textbf{Hitchin 分支}：$G$ 为分裂实群时特征簇中包含主 Fuchs 表示的连通分支，其中表示均为离散忠实；它是高阶泰希米勒空间的基本例子（见 10.1 节）。 \\
Hitchin equations & \textbf{Hitchin 方程}：黎曼曲面 $X$ 上自对偶 $G$-Yang–Mills 联络所满足的方程，引出 Higgs 丛的概念；经非阿贝尔 Hodge 对应，Higgs 丛模空间可与曲面的特征簇等同（物理上另见第 21 章的讨论）。 \\
Homology & \textbf{同调}：给拓扑空间赋予在同胚下不变的阿贝尔群的构造（单纯、胞腔、奇异同调等彼此等价），提供空间的代数不变量；其对偶概念上同调还带有环结构，光滑流形上可由闭形式模去恰当形式得到 de Rham 上同调。 \\
Homotopy & \textbf{同伦}：把两个映射连续地形变到彼此的概念，由此得到同伦群等基本的代数不变量；最简单的情形即第 1 章讨论的基本群。 \\
% [原文 p.442（印刷 p.442）]
Instanton & \textbf{瞬子}：欧氏 Yang–Mills 方程的自对偶或反自对偶经典解，通常对应非微扰的隧穿效应，在 Yang–Mills 理论的半经典分析中居核心地位（见 17.3.2 节）。 \\
Moduli space of supersymmetric vacua & \textbf{超对称真空模空间}：规范不等价真空构成的流形；在超对称理论中它常有正维数并带复凯勒、超凯勒或特殊凯勒等几何结构，还含有若干称为分支的子簇（见第 17、18 章）。 \\
Positivity & \textbf{正性}：统一 Hitchin 分支（全正矩阵）与极大表示（正定 Maslov 形式）等已知高阶泰希米勒空间例子的概念；经 $\Theta$ 正性推广后还给出 $\SO(p,q)$ 型与某些例外李群的新族（见第 12 章）。 \\
% [原文 p.443（印刷 p.443）]
Quiver & \textbf{箭图}：无 1-环与 2-环的有限有向图；它是定义簇代数簇的核心组合数据，在物理上也用来编码某些规范理论（如 $\mathcal{N}=2$ 箭图，见第 21 章）。 \\
Riemann–Hilbert correspondence & \textbf{Riemann–Hilbert 对应}：流形 $M$ 的基本群到李群 $G$ 的表示空间与 $M$ 上平坦 $G$-联络模空间（模去自同构）之间的等价；即平坦联络由其和乐完全刻画（见定理 3.30）。 \\
Seiberg–Witten theory & \textbf{Seiberg–Witten 理论}：以 Seiberg–Witten 可积系统精确描述 4d $\mathcal{N}=2$ 理论库仑分支低能物理的框架，可精确计算其特殊凯勒度量；S 类理论中该可积系统实现为相应黎曼曲面上的 Hitchin 模空间（见第 20 章）。 \\
Spectral network & \textbf{谱网络}：从属于分歧覆盖 $\tilde S\to S$ 的曲面 $S$ 上的特定图，实现 $S$ 上非阿贝尔联络与 $\tilde S$ 上阿贝尔联络之间的阿贝尔化与非阿贝尔化；在 S 类理论中它是由 BPS 轨道组成、带孤子简并度标签的网络，随角的变动可提取 BPS 谱（见第 III 部分与 22.3 节）。 \\
% [原文 p.444（印刷 p.444）]
Supersymmetry & \textbf{超对称}：联系费米子与玻色子的时空对称性，其无穷小生成元构成超庞加莱代数或其扩充；超对称场论在重整化下行为大为改善，是研究量子场论深层性质的理想框架。 \\
Teichmüller space & \textbf{泰希米勒空间}：光滑曲面 $S$ 上常曲率标记黎曼度量模去同痕（等价地，标记复结构模去同痕）的模空间；闭曲面亏格 $g\ge2$ 时同胚于 $6g-6$ 维实球（见第 6 章）。 \\
Yang–Mills theory & \textbf{Yang–Mills 理论}：把 Maxwell 理论推广到非阿贝尔规范群的规范理论，其 Yang–Mills 泛函的极值点满足 Yang–Mills 方程。 \\
\end{longtable}
\normalsize

% [Tutor Note: 「BPS state」的数学条目原文写「General representations of these algebras satisfy the inequality $M\ge|Z|$, where $M$ is a Casimir invariant and $Z$ is a (complex) central charge」，即把 $M$ 称为 Casimir 不变量；这与物理学中 $M$ 通常表示质量（Bogomol'nyi 界中的质量）不完全一致，译文照录原文措辞，读者按 $M=|Z|$ 理解即可。]

% [Tutor Note: 原文第 441 页（印刷）出现「Typical examples of higher-Teichmüller spaces」一句，其中漏掉了「rank」一词；同段前文一律用「higher-rank Teichmüller space」，据上下文此处应为「higher-rank Teichmüller spaces」。译文按正确含义处理并在本注中说明。]

% [Tutor Note: 「Spectral network」条目中的 $T^{*}_{1,0}S\supset\tilde S\to S$ 表示 $(1,0)$-形式余切丛中的谱覆盖；此处 $\tilde S$ 是谱覆盖而非万有覆叠（与符号表中 $\tilde S$ 的用法不同形同义异），须按上下文区分。]
